% Part: normal-modal-logic
% Chapter: axioms-systems
% Section: provability-properties

\documentclass[../../../include/open-logic-section]{subfiles}

\begin{document}

\olfileid[hi]{nml}{prf}{prp}

\olsection{\usetoken{S}{derivability} के गुण}

\begin{prop}\ollabel{prop:derivabilityfacts}
  मान लें $\Sigma$ कोई मोडल तंत्र और $\Gamma$ मोडल !!{formula}ों का समुच्चय
  है। निम्न गुण मान्य हैं:
  \begin{enumerate}
  \item\ollabel{prop:derivabilityfacts-monotonicity} \emph{एकदिष्टता}: यदि
    $\Gamma \Proves[\Sigma] !A$ और $\Gamma \subseteq \Delta$, तो
    $\Delta \Proves[\Sigma] !A$;
  \item\ollabel{prop:derivabilityfacts-reflexivity}
    \emph{स्वपरावर्तन}: यदि $!A \in \Gamma$, तो $\Gamma
    \Proves[\Sigma] !A$;
  \item\ollabel{prop:derivabilityfacts-cut} \emph{छेदन}: यदि $\Gamma
    \Proves[\Sigma] !A$ और $\Delta \cup \{!A\} \Proves[\Sigma] !B$
    तो $\Gamma \cup \Delta \Proves[\Sigma] !B$;
  \item\ollabel{prop:derivabilityfacts-deduction} \emph{निगमन प्रमेय}:
    $\Gamma \cup \{!B\} \Proves[\Sigma] !A$ तभी जब
    $\Gamma \Proves[\Sigma] !B \lif
      !A$;
  \item \ollabel{prop:derivabilityfacts-ruleT}% \emph{Rule T}: If
    $\Gamma \Proves[\Sigma] !A_1$ और \dots और
    $\Gamma \Proves[\Sigma] !A_n$, और
    $!A_1 \to (!A_2 \lif \cdots (!A_n \lif !B)\cdots)$ सर्वसत्यतात्मक
    दृष्टांत है, तो $\Gamma
    \Proves[\Sigma] !B$.
  \end{enumerate}
\end{prop}

प्रमाण एक सरल अभ्यास है।
\olref{prop:derivabilityfacts-ruleT} of \olref{prop:derivabilityfacts}
से, उदाहरणार्थ, यह मिलता है कि यदि $\Gamma \Proves[\Sigma] !A \lor !B$
और $\Gamma \Proves[\Sigma] \lnot !A$, तो $\Gamma \Proves[\Sigma] !B$।
इसके अतिरिक्त, आगे हम $\Gamma \cup \{!A\} \Proves[\Sigma] !B$ के स्थान पर
$\Gamma, !A \Proves[\Sigma] !B$ लिखेंगे।

\begin{defn}
  कोई समुच्चय $\Gamma$, तंत्र $\Sigma$ के सापेक्ष \emph{निगमनतः संवृत} है,
  तभी जब $\Gamma \Proves[\Sigma] !A$ से $!A \in \Gamma$ मिलता हो।
\end{defn}

\end{document}
